\documentclass{article}
\usepackage{graphicx}
\usepackage{amsmath}

\begin{document}

\section{Introduction}

Time series classification constitutes a fundamental objective within machine learning research, requiring algorithms that accurately map sequential data to discrete categories while optimising computational resource allocation. Historically, dominant methodologies have utilised intricate architectures that demand substantial processing power, thereby restricting their applicability in resource-constrained environments. Kernel-based techniques have subsequently emerged as highly efficient alternatives, leveraging randomised feature extraction to minimise training latency without sacrificing predictive performance. The original ROCKET algorithm demonstrated that untrained convolutional kernels could extract highly discriminative features in a fraction of the time required by traditional methods. The MultiROCKET framework builds upon this theoretical premise by applying multiple pooling operators and transforming the input space, thus generating a richer feature representation that fundamentally improves classification accuracy.

The selection of the MultiROCKET methodology for this reproducibility analysis is motivated by its significant structural contributions to the field and its establishment as a state-of-the-art baseline for time series evaluation. Unlike antecedent models, MultiROCKET processes both the original time series and its first-order derivative, enabling the extraction of dynamic temporal variations that might otherwise remain obscured. The architecture further refines the feature generation process by calculating the proportion of positive values in conjunction with the standard maximum pooling operation. These specific structural enhancements provide a robust theoretical justification for the reported performance improvements, as they allow the model to capture a broader spectrum of signal characteristics. Replicating the empirical results associated with these modifications provides a critical mechanism for verifying the proposed efficiency gains and ensuring the methodology remains reliable across diverse applications.

Ensuring the reproducibility of computational research is essential for maintaining scientific validity, as it allows independent researchers to verify foundational claims before adopting novel architectures. The present study is designed to systematically recreate the experimental protocols outlined in the MultiROCKET literature, providing an objective assessment of the original findings. The feature extraction pipeline and the associated linear classifier will be implemented and evaluated, facilitating a direct comparison between the reproduced metrics and the published benchmarks. Upon completing this verification phase, the underlying limitations of the current architecture will be investigated to identify potential avenues for algorithmic refinement. The integration of rigorous empirical testing with theoretical analysis will ultimately yield proposed improvements, thereby contributing to the continuous optimisation of time series classification methodologies.

\end{document}